\section{Scenarios, uncertainty and the decision layer}\label{sec:unc}

\subsection{Scenario and ensemble design}
Scenarios are choices about the world, not draws from a distribution, and we do not average them. Results are reported per pathway (SSP1-2.6, SSP2-4.5, SSP3-7.0, SSP5-8.5), with SSP5-8.5 used as a high-end stress test rather than a forecast \citep{Hausfather2020}, at 20-year windows centred on 2050, 2080 and 2100 against the fixed 1991--2020 reference. The primary ensemble is the five ISIMIP3b models; NEX-GDDP-CMIP6 supplies SSP2-4.5 and a wider model set to test whether conclusions depend on that particular five.

Two safeguards address the fact that CMIP6 contains models whose sensitivity exceeds the assessed likely range. Every headline result is indexed by global warming level (1.5, 2, 3 and 4~K above pre-industrial) as well as by calendar year, which removes model sensitivity from the argument about \emph{when}; and ensemble statistics are shown with and without any model above the assessed very-likely upper bound of equilibrium sensitivity \citep{Hausfather2022,Zelinka2020}. Performance--independence weighting \citep{Brunner2020} is a sensitivity analysis, never the default. Where overshoot runs exist (SSP5-3.4-OS), they serve as a stress test of the lag structure of Module~C: if a cohort's viability at 2100 differs between a warming-then-cooling path and a monotone path that ends at the same climate, the framework is representing legacy and hysteresis, the trajectory dependence reported for Amazon dieback under overshoot \citep{Munday2025}.

\subsection{Partitioning uncertainty}
The output $V$ is itself a probability, so the ensemble yields a distribution over probabilities. Its spread has different origins that call for different responses (Table~\ref{tab:unc}), and the framework reports them separately: a scientist can reduce knowledge uncertainty by measuring, a manager can reduce exposure to scenario uncertainty by diversifying, and neither can do anything about internal variability except plan for it.

\begin{table}[t]
\centering
\caption{Factors of the uncertainty ensemble.}\label{tab:unc}
\footnotesize
\begin{tabularx}{\linewidth}{@{}>{\raggedright\arraybackslash\bfseries\sffamily}p{3.0cm}>{\raggedright\arraybackslash}p{4.4cm}>{\raggedright\arraybackslash}p{3.0cm}>{\raggedright\arraybackslash}X@{}}
\toprule
Factor & Levels & Kind & Treatment\\
\midrule
Pathway (SSP) & 4, plus overshoot stress test & Deep (a choice) & Conditional reporting; never averaged\\
Climate model & 5 (ISIMIP3b); wider set (NEX-GDDP) & Structural & GWL indexing; with/without hot models\\
Internal variability & Years within 20-yr windows, realisations & Aleatory & Distribution of annual hazards\\
Downscaling & QDM, delta change & Structural & Both carried\\
PET formulation & 4 (\S\ref{sec:A}) & Structural & All carried\\
Soil and rooting depth & 50 draws & Epistemic & Random draws from SoilGrids quantiles and priors\\
Traits & 50 hyper-parameter $\times$ 200 individual draws & Epistemic and aleatory & Two-level Monte Carlo (\S\ref{sec:B})\\
Model class & Statistical, mechanistic, gated hybrid & Structural & All reported; hybrid is the default product\\
\bottomrule
\end{tabularx}
\end{table}

\textbf{Analysis of variance for discrete factors.} For an output $Y$ (for example regional mean $V$ or the area with $V\ge v^\ast$) evaluated on the crossed design of the discrete factors, the total variance is split by a standard ANOVA into main effects and interactions, $\mathrm{Var}(Y)=\sum_f \mathrm{Var}_f+\text{interactions}$, and the fractions $F_f=\mathrm{Var}_f/\mathrm{Var}(Y)$ are tracked as a function of lead time in the manner of \citet{Hawkins2009} and \citet{Lehner2020}. This tells us, for each horizon and region, whether the argument is about emissions, about the climate model, or about our own modelling.

\textbf{Global sensitivity for the knowledge factors.} For the continuous, epistemic factors $X=(X_1,\dots,X_k)$ (rooting depth, plant-available water, trait hyper-parameters, recalibration constants) we use variance-based indices computed on the emulator of Module~B, since the full engine is too slow for the number of evaluations needed. The first-order and total-order indices are
\begin{equation}\label{eq:sobol}
S_i=\frac{\mathrm{Var}\big[\E(Y\mid X_i)\big]}{\mathrm{Var}(Y)},\qquad S_{T_i}=\frac{\E\big[\mathrm{Var}(Y\mid X_{\sim i})\big]}{\mathrm{Var}(Y)}\ \approx\ \frac{1}{2N\,\widehat{\mathrm{Var}}(Y)}\sum_{j=1}^{N}\Big[f(A)_j-f\big(A_B^{(i)}\big)_j\Big]^2,
\end{equation}
with the Jansen-type total-order estimator recommended by \citet{Saltelli2010}, in which $A$ and $B$ are independent sample matrices and $A_B^{(i)}$ is $A$ with column $i$ taken from $B$; the cost is $N(k+2)$ emulator runs. A large $S_{T_i}-S_i$ marks interactions, which in a threshold system like this one (phase transition, stomatal closure) are expected and should not be hidden. The result is a ranked list of which measurement would reduce uncertainty most, and it sets the priorities of the trait campaign in \S\ref{sec:plan}.

\subsection{Decision layer: robust portfolios}
Ranking species by expected viability cell by cell yields a landscape planted to whichever option is best on average, and that landscape fails all at once in the scenarios where the option is not. A planting programme is a portfolio, so it is optimised as one.

Let $u$ index management units (aggregated 30~m cells, roughly 1~ha), $j$ planting options, and $\omega\in\Omega$ the members of the scenario ensemble. Let $x_{uj}\in\{0,1\}$ indicate that option $j$ is planted in unit $u$, $A_u$ its area, and $b_{uj\omega}=V_{uj\omega}\,v_j$ the benefit per hectare, viability times a value weight $v_j$ (canopy area, carbon or water services, set with stakeholders). The benefit of a plan in scenario $\omega$ is $B_\omega(x)=\sum_{u,j}A_u\,b_{uj\omega}\,x_{uj}$. The plan maximises a mixture of the mean and the mean of the worst tail,
\begin{equation}\label{eq:cvar}
\max_{x,\zeta,s}\ \ (1-\lambda)\,\frac1{|\Omega|}\sum_\omega B_\omega(x)\;+\;\lambda\Big[\zeta-\frac{1}{(1-\alpha)|\Omega|}\sum_\omega s_\omega\Big]\quad\text{s.t.}\quad s_\omega\ge\zeta-B_\omega(x),\ \ s_\omega\ge0,
\end{equation}
where the bracket is the conditional value-at-risk, the mean benefit over the worst $(1-\alpha)$ share of scenarios, in the linear form of \citet{Rockafellar2000}. The remaining constraints are
\begin{equation}\label{eq:mip}
\sum_j x_{uj}\le1,\qquad \sum_{u,j}A_u\,c_{uj}\,x_{uj}\le\mathrm{Budget},\qquad \sum_uA_u x_{uj}\le\sigma_{\max}\sum_{u,j'}A_u x_{uj'},\qquad x_{uj}\le e_{uj},
\end{equation}
one option per unit, a budget, a cap $\sigma_{\max}$ on the area share of any single option (response diversity), and an eligibility indicator $e_{uj}$ that is zero where the option is outside its applicability region or legally excluded. A per-basin cap on additional water use can be added as a further linear constraint. The problem is a mixed-integer linear programme, solved with HiGHS in the skeleton. With about 500 scenario members, $\alpha=0.8$ (the worst 100), a few thousand units and ten to fifteen options, it has of the order of $10^4$--$10^5$ binary variables, which is tractable after aggregation and, if necessary, scenario reduction.

Sweeping $\lambda$ from 0 to 1 traces the mean--CVaR frontier, and the distance between its ends is the price of robustness, expressed in units of benefit (hectares of expected surviving canopy for canopy-area weights). The optimiser exploits whatever its scenario sample happens to contain, so the reported CVaR is evaluated on a fresh, independent set of ensemble members; a large gap between the in-sample and out-of-sample values is a warning that the portfolio is over-fitted to the ensemble. Areas with a large extrapolation footprint are not silently planted or silently excluded: they are flagged, and the plan states how much of its expected benefit comes from cells outside the area of applicability.
